\chapter{Identificadores y ámbitos de definición en C-Horizon}
\label{cap:identificadores}




\section{Declaración de variables}
Para declarar una variable, se realiza de una forma muy similar a C++.
Primero de todo, se nombra el tipo de la variable y posteriormente el nombre de la misma y si es un array de $n$ dimensiones, detrás se escribe 
$ [ i_{1} ] \dots [ i_{n} ]$ donde cada $i_{j}$ representa el tamaño de la $j$-ésima dimensión. Por lo tanto, los arrays son estáticos ya que su tamaño es fijo. El nombre de una variable no puede ser cualquier cadena de caracteres, solo permitimos palabras cuyo primer carácter sea alguno de 
$["a"-"z"] \mid ["A"-"z"] \mid  "\_"$ 
y el resto de caracteres sean de
$["0"-"9"] \mid  ["a"-"z"] \mid  ["A"-"z"] \mid  "\_"$. 
Además, permitiremos declarar una variable a la misma vez que la instanciamos con un valor. Adelantamos que cada instrucción termina con \;.

\begin{algorithm}

	\Integer X \Assign \const{0} \;
	
	\Boolean Y[\const{2}] \Assign [\const{True}, \const{False}] \;
	
	\Integer Z[\const{2}][\const{2}] \Assign [[\const{0},\const{1}],[\const{2},\const{3}]] \;
\end{algorithm}



\section{Objetivos}
Descripción de los objetivos del trabajo.


\section{Plan de trabajo}
Aquí se describe el plan de trabajo a seguir para la consecución de los objetivos descritos en el apartado anterior.



\section{Explicaciones adicionales sobre el uso de esta plantilla}
Si quieres cambiar el \textbf{estilo del título} de los capítulos del documento, edita el fichero \verb|TeXiS\TeXiS_pream.tex| y comenta la línea \verb|\usepackage[Lenny]{fncychap}| para dejar el estilo básico de \LaTeX.

Si no te gusta que no haya \textbf{espacios entre párrafos} y quieres dejar un pequeño espacio en blanco, no metas saltos de línea (\verb|\\|) al final de los párrafos. En su lugar, busca el comando  \verb|\setlength{\parskip}{0.2ex}| en \verb|TeXiS\TeXiS_pream.tex| y aumenta el valor de $0.2ex$ a, por ejemplo, $1ex$.

TFGTeXiS se ha elaborado a partir de la plantilla de TeXiS\footnote{\url{http://gaia.fdi.ucm.es/research/texis/}}, creada por Marco Antonio y Pedro Pablo Gómez Martín para escribir su tesis doctoral. Para explicaciones más extensas y detalladas sobre cómo usar esta plantilla, recomendamos la lectura del documento \texttt{TeXiS-Manual-1.0.pdf} que acompaña a esta plantilla.

El siguiente texto se genera con el comando \verb|\lipsum[2-20]| que viene a continuación en el fichero .tex. El único propósito es mostrar el aspecto de las páginas usando esta plantilla. Quita este comando y, si quieres, comenta o elimina el paquete \textit{lipsum} al final de \verb|TeXiS\TeXiS_pream.tex|

\subsection{Texto de prueba}


\lipsum[2-20]